\bigskip

\textbf{Partie A :}

\medskip

Soit $h$ la fonction définie sur $\R$ par

\[h(x) = \text{e}^x - x\]

\smallskip

\begin{enumerate}
\item Déterminer les limites de $h$ en $-\infty$ et $+\infty$.
\item Étudier les variations de $h$ et dresser son tableau de variation. 
\item En déduire que :

si $a$ et $b$ sont deux réels tels que $0 < a < b$ alors $h(a) - h(b) < 0$.
\end{enumerate}

\bigskip

\textbf{Partie B :}

\medskip

Soit $f$ la fonction définie sur $\R$ par

\[f(x) = \text{e}^x\]

On note $\mathcal{C}_f$ sa courbe représentative dans un repère \Oij.

\medskip

\begin{enumerate}
\item Déterminer une équation de la tangente $T$ à $\mathcal{C}_f$ au point d'abscisse 0.
\end{enumerate}

Dans la suite de l'exercice on s'intéresse à l'écart entre $T$ et $\mathcal{C}_f$ au voisinage de $0$.

Cet écart est défini comme la différence des ordonnées des points de $T$ et $\mathcal{C}_f$ de même abscisse.

On s'intéresse aux points d'abscisse $\dfrac{1}{n}$, avec $n$ entier naturel non nul.

On considère alors la suite $\left(u_n\right)$ définie pour tout entier naturel non nul $n$ par :

\[u_n = \text{exp} \left(\dfrac{1}{n}\right) - \dfrac{1}{n} - 1\]

\begin{enumerate}[resume]
\item Déterminer la limite de la suite $\left(u_n\right)$.
\item 
	\begin{enumerate}
		\item Démontrer que, pour tout entier naturel non nul $n$,
		
\[u_{n+1} - u_n = h\left(\dfrac{1}{n + 1}\right) - h\left(\dfrac{1}{n}\right) \]

où $h$ est la fonction définie à la partie A.
		\item En déduire le sens de variation de la suite $\left(u_n\right)$.
	\end{enumerate}
\item Le tableau ci-dessous donne des valeurs approchées à $10^{-9}$ des premiers termes de la suite $\left(u_n\right)$.

\begin{center}
$\begin{array}{|l|c|}\hline
n &u_n\\ \hline
1 &\np{0,718281828}\\ \hline
2 &\np{0,148721271}\\ \hline
3 &\np{0,062279092}\\ \hline
4 &\np{0,034025417}\\ \hline
5 &\np{0,021402758}\\ \hline
6 &\np{0,014693746}\\ \hline
7 &\np{0,010707852}\\ \hline
8 &\np{0,008148453}\\ \hline
9 &\np{0,006407958}\\ \hline
10 &\np{0,005170918}\\ \hline
\end{array}$
\end{center}

Donner la plus petite valeur de l'entier naturel $n$ pour laquelle l'écart entre $T$ et $\mathcal{C}_f$ semble être inférieur à $10^{-2}$.


\end{enumerate}

\bigskip
